\icode[magenta]{N23S-MCH-SI-L4-261026-001}
\nceeID[4]{2023江苏}{10}
达·芬奇的手稿中描述了这样一个实验：一个罐子在空中沿水平直线向右做匀加速运动，沿途连续漏出沙子。若不计空气阻力，则下列图中能反映空中沙子排列的几何图形是\paren[]
	
\begin{figure}[htbp]
	\centering
	\begin{tikzpicture}[font=\scriptsize]
		\begin{scope}[scale=.8]
			\draw (-2,2) rectangle(2,-3);
			\node at(0,1.5) { 运动方向\(\longrightarrow\)};
			\draw[magenta] (0,1)--(0,-2.5);
			\node at(0,-3) [below]{\normalsize A};
		\end{scope}
		\begin{scope}[scale=.8,xshift=5cm]
			\draw (-2,2) rectangle(2,-3);
			\node at(0,1.5) { 运动方向\(\longrightarrow\)};
			\draw[magenta] (-1.5,1) parabola bend(-1.5,1) (1.5,-2.5);
			\node at(0,-3) [below]{\normalsize B};
		\end{scope}
		\begin{scope}[scale=.8,xshift=10cm]
			\draw (-2,2) rectangle(2,-3);
			\node at(0,1.5) { 运动方向\(\longrightarrow\)};
			\draw[magenta] (1.5,1) parabola bend(1.5,1) (-1.5,-2.5);
			\node at(0,-3) [below]{\normalsize C};
		\end{scope}
		\begin{scope}[scale=.8,xshift=15cm]
			\draw (-2,2) rectangle(2,-3);
			\node at(0,1.5) { 运动方向\(\longrightarrow\)};
			\draw[magenta] (1.5,1)--(-1.5,-2.5);
			\node at(0,-3) [below]{\normalsize D};
		\end{scope}
	\end{tikzpicture}
\end{figure}
\begin{answer}
	C\qquad 提示：以罐子为参考系，它是一个非惯性系
\end{answer}